%=====================================================================
\chapter{Further Reading}\label{app:resources}
%=====================================================================
\normalfigures

\section{The Books the Outline Names}

\begin{alertbox}[Two things to say about the published reading list]
The HEC 2017 and 2023 outlines give the same three books, and both editions
contain the same two defects, recorded here rather than inherited:

\begin{tightnum}
  \item \textbf{``Corman''} is a misspelling of \textbf{Cormen}, throughout
        both documents.
  \item \textbf{The 3rd edition of CLRS (2009) is cited.} The \textbf{4th
        edition (2022)} has been current since before this course was written.
        It is cited below, with the correspondence noted, because a student
        buying the recommended book today will buy the 4th.
\end{tightnum}
\end{alertbox}

\begin{conceptbox}[The three prescribed texts]
\textbf{Cormen, Leiserson, Rivest and Stein, \emph{Introduction to
Algorithms}}, 4th edition, MIT Press, 2022. (The outline cites the 3rd
edition, 2009.)

\smallskip
The standard reference, and the one this handout's structure most nearly
follows. It is a reference rather than a narrative: thorough, rigorous, and
heavy going read front to back. Use it for the chapter you are currently on.

\smallskip
\emph{What changed in the 4th edition:} machine learning and matrix
applications added; the chapter on dynamic multithreading rewritten; a new
chapter on online algorithms; Fibonacci heaps moved to the website. Chapter
numbering shifted, so a reference to ``CLRS chapter 15'' means different things
in the two editions --- this handout cites by \emph{topic}, not by number.

\medskip
\textbf{Kleinberg and Tardos, \emph{Algorithm Design}}, Pearson, 2006.

\smallskip
Better than CLRS at one thing and worth owning for it: \textbf{how to find} an
algorithm, rather than how to verify one you have been given. Its treatment of
greedy algorithms --- the exchange-argument style used in \chref{greedy} and
\chref{mst} --- is the clearest in print. Weaker on data structures.

\medskip
\textbf{Sedgewick and Wayne, \emph{Algorithms}}, 4th edition,
Addison-Wesley, 2011.

\smallskip
The implementation-first book, in Java. Its strength is exactly this handout's
measured-reality strand: it reports real running times and discusses constants
honestly. If you want to see what a professional implementation of the
algorithms in Part~II looks like, this is the book.
\end{conceptbox}

\section{Beyond the List}

\begin{conceptbox}[Worth knowing about]
\textbf{Skiena, \emph{The Algorithm Design Manual}}, 3rd edition, Springer,
2020.

\smallskip
Two halves. The first is a brisk textbook; the second is a catalogue of
problems with advice on which algorithm to reach for. The ``war stories'' ---
accounts of real problems and how they were solved or not --- are the best
answer in print to a student asking whether any of this is used.

\medskip
\textbf{Knuth, \emph{The Art of Computer Programming}}, volumes 1--4A,
Addison-Wesley.

\smallskip
The deepest analysis anywhere, in a hypothetical assembly language. Not a
course text. Volume 3 (sorting and searching) contains analyses of the
algorithms in Parts~II and III to a precision no other source attempts.

\medskip
\textbf{Dasgupta, Papadimitriou and Vazirani, \emph{Algorithms}},
McGraw-Hill, 2006.

\smallskip
Short --- about 300 pages --- and free from the authors' pages. If CLRS is too
much, this covers most of the same ground in a quarter of the length. Its
chapter on NP-completeness is a good companion to \chref{complexity}.

\medskip
\textbf{Bentley, \emph{Programming Pearls}}, 2nd edition, Addison-Wesley,
1999.

\smallskip
Old and not superseded. Column 8 is the maximum-subarray problem of
\chref{bruteforce}, including the history of how the linear algorithm was
found. The book's subject is precisely this handout's: the difference between
an idea and an implementation.
\end{conceptbox}

\section{The Original Papers}

\begin{conceptbox}[Short, readable, and worth an afternoon]
Several results in this book come from papers a student can actually read.

\rowcolors{2}{white}{lightbg}
\begin{longtable}{@{}L{3.0cm}L{6.0cm}L{5.3cm}@{}}
\toprule
\rowcolor{primary!15}
\textbf{Chapter} & \textbf{Paper} & \textbf{Note} \\
\midrule
\endhead
\chref{greedy} & Huffman, ``A Method for the Construction of
  Minimum-Redundancy Codes'', \emph{Proc.\ IRE}, 1952 &
  Three pages. Written while Huffman was a student, to avoid a final exam. \\
\chref{shortestpaths} & Dijkstra, ``A Note on Two Problems in Connexion with
  Graphs'', \emph{Numerische Mathematik}, 1959 &
  Two and a half pages, covering both MST and shortest paths. \\
\chref{strings} & Knuth, Morris and Pratt, ``Fast Pattern Matching in
  Strings'', \emph{SIAM J.\ Computing}, 1977 &
  Longer, and the prefix-function argument repays the effort. \\
\chref{complexity} & Cook, ``The Complexity of Theorem-Proving Procedures'',
  \emph{STOC}, 1971 &
  The paper that created the subject of \chref{complexity}. \\
\chref{complexity} & Karp, ``Reducibility Among Combinatorial Problems'', 1972 &
  The 21 problems. Readable, and the source of the reduction technique. \\
\chref{mst} & Tarjan, ``Efficiency of a Good But Not Linear Set Union
  Algorithm'', \emph{JACM}, 1975 &
  Where the inverse-Ackermann bound of \chref{mst} is proved. \\
\chref{structures} & Reed, ``The Height of a Random Binary Search Tree'',
  \emph{JACM}, 2003 &
  The source of the $4.311\ln n - 1.953\ln\ln n$ that \chref{structures}
  measures against. \\
\bottomrule
\end{longtable}

\smallskip
Huffman's and Dijkstra's are the two to start with. Both are shorter than the
chapters that describe them.
\end{conceptbox}

\section{Online}

\begin{conceptbox}[Used with care]
\textbf{MIT 6.006 and 6.046}, OpenCourseWare. Full lecture video from the
authors of CLRS. 6.006 maps onto Parts~I--III of this handout and 6.046 onto
Parts~IV--VI.

\medskip
\textbf{Visualgo} (\texttt{visualgo.net}). Animations of most algorithms in
this book. Genuinely useful for \chref{structures}'s trees and
\chref{mst}'s algorithms, where the shape is the point.

\medskip
\textbf{Competitive-programming references} (\texttt{cp-algorithms.com} and
similar). Excellent implementations with the constant factors already tuned.
Treat their \emph{analyses} with more caution than their code: they are
written for contest conditions, where the input distribution is known and a
worst case that cannot occur is often ignored --- which is the opposite of this
course's habit.

\medskip
\textbf{Large language models.} They will produce a plausible analysis for any
algorithm you name, including a wrong one, and they are particularly unreliable
about \emph{which} case a bound refers to and about amortised versus worst-case
cost. Use the measurement harness of \chref{project} to check anything you are
told. That is what it is for.
\end{conceptbox}

\section{A Reading Path}

\begin{conceptbox}[If you read one more thing]
\rowcolors{2}{white}{lightbg}
\begin{longtable}{@{}L{5.0cm}L{9.5cm}@{}}
\toprule
\rowcolor{primary!15}
\textbf{If you want} & \textbf{Read} \\
\midrule
\endhead
the exam answered & CLRS, the chapter matching the week \\
to design algorithms & Kleinberg \& Tardos, chapters 4 and 6 \\
to make them fast & Sedgewick \& Wayne, and Bentley \\
to know what is used in practice & Skiena, part II \\
the subject's history & the Huffman and Dijkstra papers, in one sitting \\
a shorter CLRS & Dasgupta, Papadimitriou \& Vazirani \\
to go further than this course & the deleted postgraduate syllabus in
  Appendix~F names the topics; CLRS covers flow and approximation \\
\bottomrule
\end{longtable}
\end{conceptbox}
